\section{Propagation of ordinary Hecke support between irreducible components}
\label{sec:linked-components}
\hypertarget{source-ordinary-linked-components}{}

The ordinary eigenpoints of Appendix~\ref{sec:linked-seeds}
supply the initial components. To reach the desired type-I points we pass support along the prescribed links. In rank
three one such passage can meet the reducible locus with irreducible outer plane; the
patching tests also cover this case. The resulting support statement is the input to
descent over the marked field and to the integral length
comparison. For a Hecke ideal \(J\), its image \(M_n\) in the degree-zero ring \(B_n\) controls Proposition~\ref{prop:ordinary-component-propagation} at each level: 
\[
 P\text{ is good for }J\quad\Longleftrightarrow\quad M_n\subseteq P\cap B_n.
\]


Keep the frame and the three levels of Section~\ref{sec:linked-seeds}.
Put \(d=[L^+:\Q]\), \(d_v=[L_v^+:\Q_\ell]\) for \(v\mid\ell\),
and \(D=nd+1\) in rank \(n\). We give the component comparison
with its arithmetic inputs stated explicitly. On a frame chart \(R\) write \(A=(a_{ij})\) for a frame matrix. Conjugation by the
torus gives the entry degrees from the two diagonal vectors:
\[
 r(t)A r(t)^{-1}=(t^{u_i-u_j}a_{ij})_{ij},\qquad
 \mathbf u=(1,-1)\ \text{or}\ (1,0,-1).
\]
Let \(I_{\rm up}\) and \(I_{\rm low}\) be the ideals killing,
respectively, the entries below and above the diagonal.

In
rank three the odd ideal \(I_{\rm odd}\) kills the four
entries joining the middle index to either outer index.
The complement of the triangular loci is the open subset
\[
 \operatorname{Spec}(R)\setminus
              \bigl(V(I_{\rm up})\cup V(I_{\rm low})\bigr).
\]
Within it the exceptional locus in rank three is \(V(I_{\rm odd})\). A point of type I
lies off both triangular loci and, in rank three, off the
exceptional locus. The marked normal form identifies this
with the absolutely irreducible locus. A point of type E
lies off both triangular loci but on the exceptional locus
in rank three. Its outer block is irreducible. For a Hecke ideal \(J\) let \(M_n=\Phi_n(J)B_n\).
Goodness is defined by the preimage of its vanishing locus:
\[
      P\in\operatorname{Good}(J)
                  \ \Longleftrightarrow\ M_n\subseteq P\cap B_n .
\]
Goodness specializes because contraction preserves inclusion: 
\[
 P\subseteq Q,\quad M_n\subseteq P\cap B_n\quad\Longrightarrow\quad M_n\subseteq Q\cap B_n.
\]


\paragraph{The prime conditions at the three levels.}
The support movement in Proposition~\ref{prop:marked-chain-support}
is used only with the completed local inclusions of
Appendix~\ref{sec:patching-cotangent}. To use these inclusions in the component graph we need two conclusions on the level
determinant ring. Its curve points take
values in \(k[[t]]\) with \(k\) finite of characteristic \(\ell\).
The point conditions include regular semisimple detection on the
cyclotomic kernel, the prescribed residual constituents, distinct
ordinary characters on the local test subgroups, and the local and
global restrictions of the marked level. In determinant terminology
a stable prime has absolutely irreducible determinant over an
algebraic closure of its fraction field. At a thick prime in this
application the semisimplification has an outer two-dimensional
constituent and a middle character. Write \(V_{\rm out}\) for the outer plane and \(\chi_{\rm mid}\) for the middle character. The ordinary Cayley--Hamilton algebra admits a nonsplit extension with semisimplification 
\[
 V^{\rm ss}\simeq V_{\rm out}\oplus\chi_{\rm mid},\qquad \dim V_{\rm out}=2.
\]

Let \(Q\) denote the kernel of a nice type-I curve point
in the Hecke locus. The first patching implication is
\begin{equation}\label{eq:ordinary-component-stable-test}
              P\subseteq Q\quad\Longrightarrow\quad J\subseteq P .
\end{equation}
For the second implication write \(Q_0,\ldots,Q_s\) for
the kernels of nice type-II curve points. Suppose
\(Z_i\subseteq Q_{i-1}\cap Q_i\) is thick for each
\(1\le i\le s\) and \(J\subseteq P_0\subseteq Q_0\) for
a stable \(P_0\). Then the comparison gives
\begin{equation}\label{eq:ordinary-component-thick-test}
 P\subseteq Q_s,\quad P\ \text{stable or thick}
                  \quad\Longrightarrow\quad J\subseteq P .
\end{equation}
This includes the chain with \(s=0\).
One common Taylor--Wiles system is needed for the whole list.
Section~\ref{sec:marked-patching} constructs its common tests,
gives the uniform first-order and odd obstruction bounds, and
recovers the Hecke support on the augmentation fiber. Sections~\ref{sec:patching-cotangent}
and~\ref{subsec:closed-link} give the completed comparison, including
specialization of thick primes. We verify the point conditions and
identify the ordinary modules on the three levels at the curve
construction below.

The frame constructed in Section~\ref{sec:marked-hosts}
supplies the abelian quotient bound, the local inertia generators
and the ordinary Selmer inequality \eqref{eq:fitted-relative-euler}.
The locally split nonzero residual cocycles in the outer and
adjacent weights, the bound \(C_0\), and the lower bounds on
\(d,d_v\) all apply to this frame.

\subsection{The two-dimensional corner}

We start in rank two because the outer plane will enter both rank-three comparisons.
Write \(\lambda_n\) for the ideal of \(B_n\) generated by the
uniformizer and the differences \(\psi_{v,i}(g)-1\) on inertia
with trivial residual scalar. The finite weight condition on
an ideal \(I\) in these charts is
\(\mathfrak m_{B_n}\subseteq\sqrt{I+\lambda_n}\).
It holds for \(M_n\). For the weight transport let \(q_0\)
be the residue cardinality. Raising inertial labels to
\(q_0-1\) kills their residual values.

Write \(h=q_0-1\). For units \(a,c\) with the same residue the geometric factor in \(a^h-c^h\) is a unit, because
\[
 \sum_{i=0}^{h-1}a^{h-1-i}c^i\equiv h\bar a^{h-1}\ne0 .
\] Thus equality of the powers in a
prime quotient recovers equality of the units. For the
scalar twist use its fixed inertial values in the weight
specialization of the level ring. The finite algebra over
that ring has Artinian special fiber. Its map onto the
specialized degree-zero quotient forces every prime there
to be maximal.

\begin{lemma}[Outer-corner support and the weight fibre]
\label{lem:rank-three-outer-corner}
For the indicated rank-two corner of the rank-three chart, let \(I_3\) be the ideal obtained from \(M_2\). For every type-E point
\(P\) satisfying the specified outer-block condition,
\[
                       I_3\subseteq P\cap B_3 .
\] It satisfies the finite weight-fibre condition for
\(\lambda_3\).
\end{lemma}

\begin{proof}
\proofstep{The completed rank-two comparison.}
First carry out the type-I propagation below in rank two,
where only the first patching input is needed. Quotient
\(A_3\) by the odd entries and let \(C\) be the image of
\(B_3\) there. Write \(p:B_3\to C\). The outer block
\(\sigma\) and the middle entry \(u\) in this quotient are
individually multiplicative. The outer block inherits the
transpose relation by restriction of the antidiagonal matrix
to the outer indices. It has ordinary characters \(\psi_0,\psi_2\).
To verify this integrally, at an ordinary place choose the marked
element \(g_1\) for which the middle residual value differs
from both outer values. Put
\[
 P_x=\sigma(x)-\psi_0(x),\quad
 T=\sigma(g_1)-\psi_1(g_1),\quad Q_y=\sigma(y)-\psi_2(y),
\]
where scalars multiply the identity. The determinant of \(T\)
is a unit by the marked unit condition. The rank-three ordinary
product on the outer block gives \(P_xTQ_y=0\) for all \(x,y\).

Taking \(y=g_1\) gives
\(P_xT^2=(\psi_2(g_1)-\psi_1(g_1))P_xT\).
Cayley--Hamilton in dimension two gives
\[
  \det(T)P_x =
    \bigl(\operatorname{tr}T-\psi_2(g_1)+\psi_1(g_1)\bigr)P_xT .
\]
Multiplication by \(Q_y\) yields \(P_xQ_y=0\). On the middle
entry of the triple product at \((g_1,y,g_1)\) the first and
last factors are units, giving \(u(y)=\psi_1(y)\) locally.
The corner therefore defines a map from the rank-two frame
whose degree-zero part is \(\beta:B_2\to C\). Put
\[
                        I_3=p^{-1}\bigl(\beta(M_2)C\bigr).
\]
At a type-E prime a nonzero upper and lower entry must both
lie in the outer block. Its contraction in rank two is type I,
hence good for \(M_2\). If \(P\) is the type-E prime, the rank-two contraction contains \(M_2\); taking the inverse image in the degree-zero ring gives 
\[
 I_3=p^{-1}(\beta(M_2)C)\subseteq P\cap B_3.
\]


\proofstep{The resulting weight-fibre bound on the outer corner.}
The ideal \(I_3\) also has finite weight fiber. A prime above
\(I_3+\lambda_3\) gives a prime of \(C\) whose rank-two
contraction contains \(M_2+\lambda_2\). All rank-two degree-zero
values then reduce to coefficients. Consider the middle character
\(u\) in the resulting local domain in characteristic \(\ell\).
It reduces to one, is invariant under the frame involution, and
is one on the indicated inertia kernels. Modulo each power of the
maximal ideal its continuous image is a finite \(\ell\)-group in
the principal units. The abelian generator condition of the frame
gives a fixed positive exponent \(h\) for such finite quotients
once the prescribed inertia and involution differences are killed.
Separatedness turns the finite-depth identities into an equality in the domain.

For its maximal ideal \(\mathfrak m\),
\[
 u(g)^h-1\in\bigcap_{b\ge1}\mathfrak m^b=0 .
\] Its \(\ell\)-part
can be dropped by injectivity of Frobenius. The remaining geometric
sum is a unit since \(u(g)\) reduces to one. Thus \(u=1\).
In any degree-zero monomial an odd factor vanishes and all middle
diagonal factors can be removed. The remaining monomials come
from rank two, and the middle ordinary character is one. These
generators and the coefficient reduction force the prime to be
maximal. This gives the required weight-fiber control on the
exceptional locus using only the rank-two Hecke ideal.

\end{proof}

\subsection{Curves and the component graph}

To connect components through curves, work locally at a closed point on the residual
lower triangular locus off the upper triangular locus. Use its complete pointed chart
\(R\), flat over the frame ring, and contract primes of \(R\)
to \(B_n\). Call a chart prime nice if its quotient is
one-dimensional of characteristic \(\ell\), it lies off both
triangular loci, and its distinct ordinary indices give
different characters on each distinguished local subgroup \(N_v\).
The following estimates govern the choice. In the chart the
upper triangular ideal is the unit ideal and the opposite
triangular quotient has dimension at most
\[
                   (n-1)d+2+3C_0+5 < nd .
\]
For the latter estimate cut by the uniformizer, \(1+C_0+d\)
outer diagonal values at abelian generators and \(d\) values for
each middle diagonal index.

On the triangular quotient the
diagonal characters determine the characteristic polynomials.
The cut forces the degree-zero maximal ideal and leaves chart
fiber dimension at most \(3C_0+5\). For a prime off both
triangular loci above \(I\) with finite weight fiber, coincidence
of two characters on \(N_v\) in characteristic \(\ell\) gives
\[
                    \dim(R/P)+d_v\le nd+5 .
\]
The weight count drops the \(d_v\) generators of one of the
equal labels, retaining one distinguished value, the other
generators and the coefficient uniformizer. The radical test
bounds the dimension by the number of cuts. The off-triangular
framing has at most one transcendental parameter over degree
zero. The large-frame condition gives \(d_v>5\).

\begin{proposition}[A nice specialization in the pointed chart]\label{prop:nice-ordinary-specialization}
Use a complete pointed chart at a closed point of the
residual lower triangular locus off the upper triangular locus.
Let \(P\) have quotient dimension at least \(D\), with
contraction to \(B_n\) containing an ideal \(I\) of finite
weight fibre. There exists a nice prime \(Q\) with
\[
                    P\subseteq Q,\qquad \operatorname{char}\kappa(Q)=\ell .
\] The resulting curve admits
an embedding into a complete DVR with finite residue field.
\end{proposition}

\begin{proof}
\proofstep{Producing a curve with the required specializations.}
Let \(P\) have quotient dimension at least \(D\), with contraction
containing such an \(I\). Imposing the uniformizer leaves a
component of dimension at least \(nd\), by the principal ideal
theorem and catenarity. The bounds just proved keep its generic
point off the opposite triangular locus and all local coincidence
loci. In that component choose a product \(f\) of one nonvanishing
element from each bad ideal, times a nonzero element of the maximal
ideal.

Let \(A\) be the local domain of the chosen component, and write \(R\twoheadrightarrow A\) for the quotient of the pointed chart. Lemma~\ref{lem:local-curve-avoidance} applied to \(0\ne f\in\mathfrak m_A\) gives a prime \(\mathfrak q\) avoiding \(f\) with one-dimensional quotient. Pull it back to \(Q\) in \(R\). Then
\[
 P+(\varpi)\subseteq Q,\qquad R/Q\simeq A/\mathfrak q,
 \qquad f\notin\mathfrak q .
\]
The first inclusion gives the specialization and its residue characteristic. The last condition keeps its generic point away from every excluded locus, since the factor of \(f\) chosen in the corresponding ideal survives there. The complete local domain of the
resulting curve embeds into a complete DVR \(k[[t]]\) with
\(k\) finite. A parameter gives a finite extension of a
one-variable power series ring over the coefficient field. Normalizing in the finite fraction field
extension gives the DVR.

\end{proof}

\paragraph{The tests on the three level rings.}
\label{par:fitted-arc}
\hypertarget{source-ordinary-marked-points}{}
The curves lie in a frame. Their specializations at the three levels must pass the
tests of Section~\ref{sec:marked-patching}. In this
paragraph \(P=\bar\Psi\) denotes the character of the fitted field
construction in Section~\ref{sec:marked-hosts}. The residual characters
on the CM group at the level are
\(\bar\kappa_n P^{w_i}\), with \(w=(1,-1)\) or \((1,0,-1)\).
The point on the level is the pullback of the chart representation
by the surjection to the frame group, multiplied by \(\kappa_n\).
Its ordinary characters are multiplied by the same scalar. Thus the
specialized determinant and ordered characters evaluate by the map
from the actual level ring to the chart followed by the map to the
curve DVR. In particular its kernel is the required contracted prime.
The finite scalar and the cyclotomic twist have the prescribed
polarization, as in Section~\ref{sec:linked-seeds}.

\begin{auxlemma}[Multiplicity of residue characters]\label{lem:residue-character-multiplicities}
Let \(k\) be a field of characteristic \(\ell>0\), and let \(\chi\) range over a finite set of distinct characters \(H\to k^\times\). If two lists of nonnegative integer multiplicities satisfy
\[
 \prod_\chi (T-\chi(h))^{a_\chi}
       =\prod_\chi (T-\chi(h))^{b_\chi}\quad\text{in }k[T]
       \qquad(h\in H),
\]
then \(a_\chi=b_\chi\) for every \(\chi\).
\end{auxlemma}
\begin{proof}
The coefficient below the leading term gives \(\sum_\chi (a_\chi-b_\chi)\chi=0\) as functions on \(H\), with the differences read in \(k\). Distinct characters are linearly independent, so \(a_\chi-b_\chi\in\ell\mathbb Z\). Cancel the \(\min(a_\chi,b_\chi)\) common factors for each \(\chi\). The remaining multiplicities on both sides are multiples of \(\ell\), giving
\[
 \left(\prod_\chi(T-\chi(h))^{(a_\chi-b_\chi)_+/\ell}\right)^\ell
 =
 \left(\prod_\chi(T-\chi(h))^{(b_\chi-a_\chi)_+/\ell}\right)^\ell ,
 \qquad u_+=\max(u,0).
\]
Frobenius is injective on \(k[T]\), so the two inner polynomials agree at each \(h\). The remaining total degree strictly decreases at every iteration that starts in positive degree. It reaches zero, proving equality of all original multiplicities.
\end{proof}

\begin{proposition}[Local labels and spans at a fitted point]
\label{prop:fitted-points-three-tests}
At the indicated points of the three level rings the ordered labels
remain distinct on the specified local groups after restriction
to an open subgroup. Write \(H\) for the full cyclotomic kernel on the positive component.
Over an algebraic closure of the coefficient field the two block spans are
\[
 \operatorname{span}r(H)=
 \begin{cases}
   \operatorname{End}(V),&\text{type I},\\
   \operatorname{End}(V_{\rm out})\oplus\operatorname{End}(V_{\rm mid}),&\text{type E},
 \end{cases}
\]
where \(V_{\rm out}\) is the outer plane and \(V_{\rm mid}\) the middle line. In the
open-coset test the element with separate eigenvalues is chosen
in that coset.
\end{proposition}

\begin{proof}
\proofstep{Testing the ordered local labels.}
At each \(v\mid\ell\), on the preimage of the distinguished local
subgroup \(N_v\) the residual character ratios are one and the full
cyclotomic character is one. The former follows from \(P|_{N_v}=1\),
and the latter is part of the definition of \(N_v\).
The nice curve avoids every local coincidence ideal, so each
distinct pair of labels has a nonidentity ratio somewhere on this
subgroup. This persists under the scalar twist. The group
\(1+t k[[t]]\) is torsion-free in characteristic \(\ell\):
raising to \(\ell^a\) raises the degree of the leading nonconstant
term by that factor, and raising to a power prime to \(\ell\)
multiplies its coefficient by a unit.

For a ratio \(\alpha\) nontrivial on an element \(n\) in
the preimage of \(N_v\), choose \(m>0\) with \(n^m\) in any
prescribed open subgroup of the global group. Then
\[
    \alpha(n^m)=\alpha(n)^m\ne1,\qquad
    \epsilon(n^m)=1 .
\]
Thus the same pair of labels is distinct on the restricted
cyclotomic kernel. We can choose a single element there on which
all ratios have infinite order. To see this, successively multiply by
powers of elements on which the individual ratios are nontrivial. A previously nontrivial
ratio can become one for at most one exponent if it is nontrivial
on the new element, and otherwise stays unchanged. The new ratio
can also be made nontrivial by avoiding one exponent. This supplies
an element of the local group in the chosen open subgroup and the
full cyclotomic kernel with pairwise distinct eigenvalues at
\emph{every} positive power. At each marked place the ordering must pass both tests simultaneously:
\[
 \psi_i\ne\psi_j,\qquad \psi_i/\psi_j\ne\epsilon\qquad(i<j).
\]

\proofstep{Full spans on the cyclotomic kernel.}
We next check absolute irreducibility on small subgroups of the
cyclotomic kernel. Write \(G\) for the level CM group and
\(H=\ker(\epsilon|_G)\). Work over an algebraic closure of the
fraction field of the curve. In type I the marked representation is
irreducible on the frame group, hence on \(G\) after pullback and
scalar twist. In type E apply this assertion to its outer plane,
irreducible since neither outer triangular ideal vanishes. Its two
residual character labels are the outer labels above. Indeed on the
residual triangular chart the diagonal entries give those characters
and the product of the two opposite outer entries is zero.

On the outer plane the ordinary labels are obtained by removing the middle character with multiplicity one. If \(S\) denotes the multiset for the full flag, then
\[
 S=S_{\rm outer}\sqcup\{\chi_{\rm middle}\}.
\] One way to see the deletion
is to take a local element with three distinct eigenvalues. Both
summands and all flag steps are sums of its eigenlines, so their
intersections give the induced flags.

In either irreducible block the dimension \(b\) is a prime, two or
three. Let \(U\) be any open subgroup of \(G\) and take
\(H'=H\cap\operatorname{core}_G(U)\), normal in \(G\).
If \(W_0\) is a simple \(H'\)-submodule, every \(gW_0\)
is again simple for \(H'\) by normality. Their sum is a
nonzero \(G\)-submodule, hence the full irreducible block.
This proves semisimplicity. Irreducibility for \(G\) makes
the action on \(H'\)-isotypic pieces transitive. If \(t\)
counts the pieces, \(d_1\) is the dimension of a simple
piece's constituent and \(m_1\) its multiplicity, then
\[
                             b=t\,d_1\,m_1 .
\]
Since \(b\) is prime, reducibility forces either \(t=d_1=1\)
(scalar action) or \(d_1=m_1=1\) and \(t=b\) (distinct
character lines permuted transitively).
The separating local element in \(H'\) excludes scalars.

In the
second case let \(G_1\) be the kernel of the permutation action.
On the intersection of the local decomposition group with \(H'\),
the characters of the lines agree up to permutation with the
restrictions of the ordinary labels. This follows by comparing
the semisimplification of the flag with the direct sum of the lines.
Write \(D\) for the local decomposition group, and \(\lambda_i\) for the distinct characters on \(D\cap H'\). Each is the restriction of a character \(\widetilde\lambda_i\) of \(D\); hence for \(d\in D\) and \(h\in D\cap H'\),
\[
 \lambda_i(dhd^{-1})=\widetilde\lambda_i(d)\lambda_i(h)\widetilde\lambda_i(d)^{-1}
 =\lambda_i(h).
\]
The corresponding eigenspaces are distinct lines. Every \(d\) preserves each line, so \(D\subseteq G_1\).

On \(G_1\) the line characters are units of the valuation ring of the
algebraic closure: their values and inverses are roots of monic
characteristic polynomials over the curve DVR. Reducing gives
characters whose product characteristic polynomial at each element
is that of the prescribed global residual labels restricted to
\(G_1\). The values of the reduced line characters belong to the finite image of the residual diagonal characters. The two lists consequently factor through a common finite abelian quotient of order prime to \(\ell\). Applying Lemma~\ref{lem:residue-character-multiplicities} on \(G_1\) identifies their multisets with their full integer multiplicities. Write \(\chi\) for a reduced character on \(G_1\) and \(\tilde\chi\) for a global extension. Then conjugation fixes \(\chi\), since
\[
 \chi(ghg^{-1})=\tilde\chi(g)\chi(h)\tilde\chi(g)^{-1}=\chi(h)\qquad(h\in G_1).
\] Transitivity of the original lines forces their
reductions to agree on \(G_1\), including the local group.

The local
mark with \(P(g_*)^2\ne1\) distinguishes the outer residual labels
there. This contradiction proves irreducibility on \(H'\), hence on
\(U\cap H\). The local mark exists at each place by conjugating
the mark of \eqref{eq:tag-h1} in the abelian fitted field and using the
global character \(P\). Equivalently one place suffices for this
contradiction. In type E the two simple \(H\)-constituents are inequivalent
because their dimensions are different. The semisimple image
algebra of the group ring acts as the product of their full
endomorphism algebras by density. In the two cases the spans
over the algebraic coefficient field are therefore
\[
  \operatorname{span}r(H)=
   \begin{cases}
      \operatorname{End}(V),&\text{type I},\\
      \operatorname{End}(V_{\rm out})\oplus\operatorname{End}(V_{\rm mid}),
          &\text{type E}.
   \end{cases}
\]
The equality follows over the original coefficient field by
extension of scalars, since the absolute dimensions are the
displayed block dimensions.

\paragraph{Regular elements in each open coset.} Both block patterns now require the same density statement, with the group element
constrained to a given open coset. In an open coset meeting \(H\)
choose \(s\in H\) and an open subgroup \(U\) with
\(s(U\cap H)\) inside the coset. Let \(C\) be the Zariski closure
of the image of \(U\cap H\) and \(C^\circ\) its identity component
over the algebraic closure. The preimage of \(C^\circ\) is open in
\(U\cap H\). Indeed each of finitely many components is closed for
the adic topology on matrix entries: individual polynomial equations
over the algebraic closure already have coefficients in a finite
extension. Choose an open subgroup \(V\subset U\) of \(G\) small enough that 
\[
 V\cap H\subseteq r^{-1}(C^\circ),\qquad r(V\cap H)\subseteq C^\circ.
\]
 The irreducible blocks just proved
are irreducible for \(C^\circ\).

\proofstep{A diagonal torus in the connected image.}
Take the separating local element in \(V\cap H\) and diagonalize its
matrix. The identity component of the closure of its powers is a
torus inside the diagonal torus of \(C^\circ\). Let \(x\) be the matrix of this element and let \(T\) be the identity component of the closure of its powers. For distinct eigencharacters \(\lambda_i,\lambda_j\), coincidence on \(T\) would yield 
\[
 (\lambda_i/\lambda_j)(x)^m=1,\qquad m=[\overline{\langle x\rangle}^{\rm Zar}:T],
\]
 contrary to the separating property. In the rational cocharacter space the equal-coordinate loci are finitely many proper hyperplanes. Choose a rational point in their complement and clear denominators to obtain the integral cocharacter
\[
 d=(d_i),\qquad \prod_{i<j}(d_i-d_j)\ne0 .
\]

In this eigenbasis there is \(h\in C^\circ\) such that all principal
minors of \(A=r(s)h\) are nonzero. For one block of dimension at
most three the first and the next-to-last exterior powers are
irreducible: the latter is the dual of the first tensored by the
determinant. The other nonempty exterior power is the determinant
itself. For each principal minor, the orbit of its basis wedge
therefore spans the corresponding exterior power.

For a wedge \(e_I\) and its coordinate functional \(e_I^*\), the principal minor equals 
\[
 \det (r(s)h)_{I,I}=e_I^*\!\left(\bigl(\textstyle\bigwedge^{|I|}(r(s)h)\bigr)e_I\right).
\]
 The orbit spans and the functional composed with \(\bigwedge^{|I|}r(s)\) is nonzero, so this minor is nonzero somewhere on \(C^\circ\).
For the outer plane plus line, a principal minor is the product of
a principal minor on the plane and, when the line is included, its
nonzero scalar. The same argument on the plane works. The simultaneous
choice follows from irreducibility of the connected algebraic group.

\proofstep{A separable characteristic polynomial in the desired coset.}
Consider \(A\operatorname{diag}(z^{d_i})\) over the Laurent
series field in \(z\). Its characteristic coefficient at \(X^{n-j}\)
is
\[
 (-1)^j\sum_{|I|=j}\det A_{I,I}\,z^{\sum_{i\in I} d_i}.
\]
The subset of the \(j\) smallest exponents uniquely minimizes
the order and its principal minor is nonzero. Number the ordered
exponents \(d_{(1)}<\cdots<d_{(n)}\). The consecutive coefficient
orders, read from the leading term, differ by the successive
\(d_{(i)}\). The Newton polygon therefore has \(n\) segments of
horizontal length one with distinct slopes. Its roots
over the algebraic closure have distinct valuations. Hence the
discriminant on \(r(s)C^\circ\) is a nonzero regular function, since
its pullback along this cocharacter translate is nonzero.
The image of \(U\cap H\) is Zariski dense in \(C\), so some actual
element of \(s(U\cap H)\) has nonzero discriminant. This proves the
required density in the profinite topology, including cosets outside
the identity component.

\end{proof}

\paragraph{Integral normalization and the distinguished place.}
To apply the tests at the distinguished place we need integral representatives of the
blocks. After a finite coefficient extension use coordinates on the blocks
defined over the fraction field of the new DVR. In type E their
dimensions differ, so polarization pairs each block with itself:
a cross pairing would intertwine inequivalent simple modules after
conjugation and duality. The form on each block is nondegenerate
and symmetric.

In a block \(V_i\) over a coefficient extension \(E'/E\), intersect the extended lattice with that subspace: 
\[
 L_i=(L\otimes_{O_E}O_{E'})\cap V_i,\qquad r(g)L_i=L_i.
\]
 Here \(L\) is the original stable lattice and \(O_{E'}\) the extended coefficient DVR. The dual is stable as well. Indeed, the polarized action transports integral pairings to integral pairings, since its multiplier is a unit and conjugation preserves the group. Applying the same argument to the inverse gives, for the dual lattice \(\Lambda^\vee\),
\[
 r(g)\Lambda^\vee\subseteq\Lambda^\vee,\quad r(g^{-1})\Lambda^\vee\subseteq\Lambda^\vee\quad\Longrightarrow\quad r(g)\Lambda^\vee=\Lambda^\vee .
\] Diagonalize the Gram matrix over the DVR,
allowing fractional entries. One may first scale to an integral
matrix and choose a primitive vector with self-pairing of smallest
valuation among the entries. When a mixed entry has that valuation
and both diagonal entries have greater valuation, take the sum of
the two basis vectors, using \(2\) invertible. Split off this
vector by integral row and column operations and continue.

Write the diagonal entries as \(t^{a_i}u_i\) with \(u_i\) units.
In this basis preservation of the lattice and its dual requires
\[
       v_t(r_{ij})\ \ge\ \max(0,a_j-a_i).
\]
After adjoining \(\sqrt t\), use the basis \(e'_i=t^{-a_i/2}e_i\).
In the extended valuation with \(v_t(t)=1\), the new entries satisfy
\[
  v_t(r'_{ij})=v_t(r_{ij})+\tfrac12(a_i-a_j)
                 \ \ge\ \tfrac12|a_i-a_j|.
\]
The same calculation applies to inverse matrices, so the new lattice
is stable. Its Gram matrix is \(\operatorname{diag}(u_i)\).
Adjoin square roots of the reductions of the \(u_i\) by a finite
residue extension and lift them by Hensel's lemma. Put \(e_i=u_i^{-1/2}e'_i\). Since the chosen square roots are units, this preserves the lattice and normalizes each diagonal Gram entry:
\[
 (u_i^{-1/2})^2u_i=1 .
\]

In each ordinary flag a local matrix with distinct
eigenvalues splits its steps into eigenlines over the fraction
field. These eigenlines lie in the separate blocks.
Intersect the resulting flags in each block with its lattice.
The steps are saturated, and successive free quotients over the
DVR give an integral basis compatible with their order. Interleave
the block bases in the prescribed local order. The original and new
coordinates are related over a finite extension by fixed matrices,
so adic continuity, the two character inequalities, and the span
and separability tests persist. The coefficient maps are injective
and local, hence keep the evaluation kernels and determinant
identities. This supplies the normalized points used for the
simultaneous prime choice.

\begin{lemma}[The marked support tests on the chosen curves]
\label{lem:fitted-marked-support-test}
In the integral normalization the distinguished-place test agrees
with the auxiliary level maps of the three rings. Let \(b\) be the trace Gram minor nonzero at the type-I point \(Q\).
The same minor detects the full matrix span at its generizations since
\[
           P\subseteq Q,\quad b\notin Q\quad\Longrightarrow\quad b\notin P.
\]
The support test therefore applies at the indicated type-I generizations;
the two-block and augmentation comparisons give support at the
indicated stable or thick type-E target from the stable start.
\end{lemma}

\begin{proof}
\proofstep{The distinguished-place vanishing test.}
There is a split distinguished place for which the twisted dual
test holds for \emph{all} these curves. On the field \eqref{eq:tag-h2} choose
a CM element with \(P=1\) and \(\bar\epsilon\ne1\). Explicitly
multiply complex conjugation by the element with \(P=\bar\epsilon=1\),
\(\eta=-1\) in Section~\ref{sec:marked-hosts}. If necessary multiply
by inertia at the new prime ramified in \(\Q(\sqrt{Dp})\) to enter
\(H_0\); the three characters \(P,\bar\epsilon,\eta\) are unramified
there. Raise to a common \(\ell\)-power to enter the three
\(\ell\)-power kernels of \eqref{eq:tag-h2}.

If \(g\) is the adjusted element, the common \(\ell\)-power preserves its cyclotomic residue because \(\ell\) is odd: 
\[
 \bar\epsilon(g^{\ell^a})=\bar\epsilon(g)^{\ell^a}=(-1)^{\ell^a}=-1.
\]
 Chebotarev over \(L^+\) for the finite quotient recording
\(P,\bar\epsilon,\eta\) gives places \(v_0\) split in \(L\) with
such Frobenius, away from the cutoff and the discriminant. One can require residue degree one over \(\mathbb Q\). Indeed the higher-degree primes of a fixed number field \(F\) contribute a bounded sum as \(s\downarrow1\), since
\[
 \sum_{\mathfrak p:\ f(\mathfrak p/p)>1}(N\mathfrak p)^{-s}\le [F:\mathbb Q]\sum_p p^{-2s}=O_F(1).
\] Choose the second omitted place away from
the first. At each chosen decomposition group all the twisted
residual characters have the common value \(\bar\kappa_n\).

For a Frobenius matrix \(T\) at an arc and cyclotomic value \(e_0\),
the reductions of \(T\) and \(e_0 T\) have single, distinct,
nonzero eigenvalues. Their characteristic polynomials have unit
resultant over the curve DVR. Let \(P,Q\) be the characteristic polynomials of
\(e_0 T,T\) over the curve DVR. Since the resultant is a unit, the two polynomials generate the unit ideal, giving a B\'ezout identity over the coefficient ring:
\[
 AP+BQ=1 .
\]
If \(e_0TY=YT\), then \(f(e_0T)Y=Yf(T)\) for every
polynomial \(f\). Applying the Bézout relation on the
left gives
\[
  Y=A(e_0T)P(e_0T)Y+B(e_0T)YQ(T)=0
\]
by Cayley--Hamilton. This only needs the reduced
characteristic polynomials, so it covers nonsemisimple
reduction. This gives the distinguished-place vanishing test and the nonidentity cyclotomic scalar used in the paired detector.

At each of the three levels take the finite set consisting of all
places above \(\ell\), the tame places above the cutoff away from
\(\ell\), and its own distinguished \(v_0\). All split in the CM
extension, which is unramified at finite places. The base Galois
level has exactly this set of allowed ramification. At the
distinguished place use the first principal congruence subgroup,
at the fixed tame places the coefficient type for the residual
inertia characters, and at \(\ell\) the ordered ordinary torus
character, with the twists of Section~\ref{sec:linked-seeds}.

If \(u\) has order dividing \(\ell^a\), its value under a local character \(\theta\) into the curve DVR satisfies
\[
 (\theta(u)-1)^{\ell^a}=\theta(u)^{\ell^a}-1=0.
\]
The DVR is reduced, so \(\theta(u)=1\). Under local reciprocity at a split place \(w\mid\ell\), the free pro-\(\ell\) unit rank is \([L_w:\Q_\ell]\); a uniformizer supplies the independent unramified generator.
The base evaluation kernel on the empty-auxiliary Hecke action is
precisely \(J_2,J_3\) or \(J_0\) according to the chosen ring.

\proofstep{Compatibility with the auxiliary level maps.}
For each selected system of \(q\) further places use the two maps
of Section~\ref{par:ordinary-level} for that same ordinary module.
Their composites are the same split-root factor, the tame
augmentation is given by the trace, and the relation bound there is
\(n!\) at every auxiliary depth. The proofs use rank \(n\) compact
and Iwahori fixed vectors, regular ordinary specialization and the
fixed tame type. At the two marked levels the relation bound \(n!\) reads
\[
 2!=2\quad(n=2),\qquad 3!=6\quad(n=3)
\] Choose one split completion \(w\) above each \(v\mid\ell\) of \(L^+\), so \(L_w\simeq L^+_v\). The \(n\) ordinary unit characters therefore give
\[
 n\sum_{v\mid\ell}[L^+_v:\Q_\ell]=n[L^+:\Q]
 \quad(n=2\text{ or }3).
\]
The two rank-three actions are compared to their own base ideals before transporting either to the common frame.

The ordinary deformation and the ordered characters of each curve
are defined by the same pullback and scalar as the map \(\Phi_n\).
Let \(Q\) be the type-I curve kernel and \(b\) its nonzero trace Gram minor. Every generization \(P\) satisfies 
\[
 P\subseteq Q,\quad b\notin Q\quad\Longrightarrow\quad b\notin P.
\]
 Starting at that point in the Hecke support,
the one-block local comparison and augmentation put every such
generization in the support.

At a chain of type-E points the two
blocks above give the hypotheses of the two-block local comparison.
Sections~\ref{subsec:closed-link} and~\ref{subsec:marked-augmentation}
then give support on the stable or thick target from a stable
start in the support. These are the prime conditions at the three
marked ideals under the ordinary level construction of this
supplement. Constants for curve normalization, cohomology and the
number of tests may depend on the fitted field, level and finite
list of curves. The patching conclusion is ideal containment and
introduces no such factor in the length estimate on the original field. \label{cor:ordinary-marked-points-1}
\end{proof}


\paragraph{Contraction to a level.} Now contract to a level. A type-I generization gives
a stable determinant, since the representation
is absolutely irreducible and the level group maps onto the
frame group. For a thick generization use the nonsplit
ordinary realization of its block determinant. Its group algebra
map kills the Cayley--Hamilton relations by the characteristic
polynomial. The ordered products, before and after the involution,
map to scalar multiples of the ordinary and reversed ordinary
products of the extension. The tame inertia relations follow
from the marked prime-to-\(\ell\) power test. Adic continuity
kills the closure of the ideal of relations. Thus the nonsplit
realization factors through the ordinary Cayley--Hamilton algebra,
while its determinant is reducible. Contraction transports
both prime conditions, including the links that must be thick.

\begin{construction}
\begin{tikzpicture}[>=Stealth,every node/.style={align=center,inner sep=2pt}]
\node (two) at (0,0) {Rank-two comparison on \(M_2\)};
\node (input) at (0,-.85) {Outer-corner support};
\node (three) at (-2.3,-1.9) {\(M_3\) and its quotient};
\node (zero) at (2.3,-1.9) {\(M_0\) and its quotient};
\draw[->] (two)--(input);
\draw[->] (input)--(three);
\draw[->] (input)--(zero);
\end{tikzpicture}
\caption{Propagation uses the rank-two comparison first. Its outer-corner input is applied to two different rank-three support problems.}\label{fig:ordinary-three-levels}
\end{construction}

\paragraph{Which components can link curves.}
Which components do these links reach from the seed?

\begin{proposition}[Propagation from a seed across ordinary components]
\label{prop:ordinary-component-propagation}
In the selected ordinary frame, given the seed of Appendix~\ref{sec:linked-seeds} of the required dimension
and the completed comparisons
\eqref{eq:ordinary-component-stable-test} and~\eqref{eq:ordinary-component-thick-test},
each type-I prime \(P\) is good for its corresponding level ideal, meaning 
\[
 M_n\subseteq P\cap B_n.
\]

The application to \(M_2\) supplies the outer-corner input for the
applications to \(M_3\) and \(M_0\), which keep their separate
quotient maps.
\end{proposition}

\begin{proof}
\proofstep{Relative dimension and the presentation of the chart.}
The ungraded chart has a presentation over a complete DVR in
\(N\) variables with \(N-D\) relations, padding with zeros if
necessary. At a rank-three point on the exceptional locus, use
conjugation by the sign on the middle coordinate to grade the
presentation by parity. There are \(a\) even and \(g\) odd
variables, at most \(a-D\) even and \(g\) odd relations. The relation count uses the relative obstruction group.
For a small extension with kernel \(I\), let \(c(g,h)\) be
the \(I\)-valued error in multiplying two matrix lifts.
For each ordinary place let \(b_v(g)\) be the error in the
lifted local flag, read in the adjoint quotient by
\(\mathfrak b_v\). With the cone differential
\(d(c,b)=(dc,\operatorname{res}(c)-db)\), associativity
and the local multiplication law give
\[
      dc=0,\qquad db_v=\operatorname{res}_v(c).
\]

Write \(\alpha\) for the one-cochain changing the matrix
lifts and \(\beta_v\) for the zero-cochain changing the flag at
\(v\), in the cone convention just specified. The new errors are
\[
 c'=c+d\alpha,\qquad
 b'_v=b_v+\operatorname{res}_v(\alpha)-d\beta_v .
\]
Thus a zero obstruction class permits both errors to be removed
in the same lift. The dual
of the minimal relation space injects into this obstruction
space. Indeed if a nonzero functional gave zero obstruction,
push out along it and lift a section to the power series variables.
Two lifts of the variables then differ in the square-zero
kernel killed by the maximal ideal. Every minimal relation has
zero relative linear part modulo that ideal, so its value is
unchanged. A relation on which the functional is nonzero would
have both zero and nonzero value. Since \(2\) is invertible, an involution \(s\) admits the complementary projections
\[
 e_+=(1+s)/2,\qquad e_-=(1-s)/2,\qquad e_+e_-=0.
\]
Applying the argument to the images gives both eigenspace assertions.

In the ordinary tangent and Euler characteristic formulas of
the marked frame, the finite-place contribution is
\(\sum_{v\mid\ell}d_v\dim\mathfrak b_v\), from which one subtracts the
sum of the dimensions fixed by real conjugations. Here
\(\dim\mathfrak b_v=n(n+1)/2\). Multiplication by the
polarization matrix identifies the fixed matrices at a real
place with skew matrices, of dimension \(n(n-1)/2\).

Since \(\sum_{v\mid\ell}d_v=d\) and \(L^+\) has \(d\) real
places, counting the extra frame direction gives
\[
 d\frac{n(n+1)}2-d\frac{n(n-1)}2+1=nd+1=D .
\]
In rank three the Borel dimensions are \(4\) even and \(2\) odd:
the even part consists of the three diagonals and the outer root.
The real skew dimensions are \(1\) even and \(2\) odd, from the
outer pair and the two adjacent pairs respectively. The extra
frame direction is even. Thus the two lower bounds for relative
tangent dimension minus obstruction dimension are
\[
        4d-d+1=D,\qquad 2d-2d=0.
\]

\proofstep{Connectedness and the odd incidence equations.}
In rank two, or outside the exceptional locus of rank three,
Grothendieck connectedness applied to the ungraded presentation
joins the minimal components in a graph with intersections
of dimension at least \(D\). Each component has dimension at
least \(D+1\). At the remaining rank-three centers write \(y\)
for the odd variables. If there are odd variables, pad the odd
relations to a square system \(H y=0\). Each odd relation lies
in the ideal \((y)\) by parity, and minimality permits the
entries of \(H\) to be taken in the maximal ideal. Impose
\(\det H=0\) in the chart. Every type-I prime contains the extra equation. Indeed modulo such a prime at least one entry of \(y\) survives, whereas
\[
 0=\operatorname{adj}(H)Hy=(\det H)y.
\]
The quotient is a domain, so \(\det H\) vanishes there. The odd dimension on this pointed chart is positive in the exceptional case.

To see the same intersection bound, use an independent
\(g\times g\) matrix \(X\) and vector \(z\). The incidence
locus \(Xz=0,\det X=0\) has irreducible support of codimension
\(g\). On \(z_i\ne0\) solve for the \(i\)-th column, with
all other coordinates free, and these charts overlap. At \(z=0\), choose a nonzero vector \(v\in\ker X\) for the singular residue matrix. The line is a specialization from the open locus, since
\[
 X(tv)=0,\qquad tv\ne0\ (t\ne0),\qquad (tv)|_{t=0}=0 .
\] The homogeneous prime radical gives a domain after
completion: the initial homogeneous pieces of two nonzero
elements still have nonzero product. Keep the \(a\) even and
\(g\) odd variables of the original chart separate from the
\(g^2+g\) new variables \((X,z)\). The new incidence equations
have codimension \(g\) in the new variables. Thus the ambient
irreducible complete local germ has dimension
\[
             1+a+g+(g^2+g)-g=1+a+g+g^2.
\]

In this germ impose \(X-H=0\), \(z-y=0\) and the at most
\(a-D\) even relations. Their number is at most
\(g^2+g+a-D\), so each component has dimension at least
\(D+1\). The complete local connectedness theorem with the
same equation count gives connectedness in dimension \(D\).
\cite{stacks-connectedness}

\proofstep{The nonsplit extension at an exceptional branch.}
The minimal primes of this smaller locus are off both triangular
loci by the bounds above. An exceptional prime there has \(y=0\)
and \(H\) singular. Choose a nonzero vector in its kernel
over the fraction field and clear denominators. At \(y=0\) an odd variation annihilates the even relations. For the odd ones the linearization is
\[
 d(Hy)=(dH)y+H\,dy=H\,dy=0 .
\] This gives a derivation
of the original chart, zero on even entries and ordinary characters.
If \(r_{\rm out},r_{\rm mid}\) are the two retained blocks,
let \(u(g)\) be the derivative of the column from the middle
line to the outer block. The even derivatives vanish, so the
linearized multiplication law in that column is
\[
   u(gh)=r_{\rm out}(g)u(h)+u(g)r_{\rm mid}(h).
\]
Consequently \(\left(\begin{smallmatrix}r_{\rm out}(g)&u(g)\\
0&r_{\rm mid}(g)\end{smallmatrix}\right)\) is a representation with
the original block determinant.

The same linearization of the
two ordered ordinary products gives the prescribed local tests. A splitting would give an
invariant eigenvector with middle coordinate 1. Write a splitting vector as \((u,1)\) in the outer-plane
and middle-line decomposition. Its equations at the mark
\(g_*\) and its polarized mate \(g_*^*\) have a common
off-diagonal column and the same middle value, while the
two distinct outer eigenvalues are exchanged. Write \(D_*,D_*'\) for the two outer diagonal matrices, \(c\) for
the common column and \(m\) for the common middle value. The two splitting
equations give
\[
 D_*u+c=mu=D_*'u+c
 \quad\Longrightarrow\quad (D_*-D_*')u=0 .
\]
Both diagonal entries of the difference are nonzero, forcing \(u=0\). The splitting equation at all
\(g\) then says the differentiated column vanishes. Polarization also kills the opposite differentiated
row, contradicting the nonzero odd derivation. This proves
thickness at the exceptional branches used by the graph.

\proofstep{Propagation from a seed.}
Suppose a chart has a good type-I prime with quotient dimension
at least \(D\). Write \(P\) for this good prime. Using \(I=M_n\), choose the minimal prime \(P_{\min}\) and the nice specialization \(Q\) so that
\[
 P_{\min}\subseteq P\subseteq Q,\qquad \dim(\text{chart}/P)\ge D .
\] The minimal prime is
type I by generization, and is a vertex of the component graph,
since every type-I prime contains the additional chart equation
when one is imposed. If the nice point is type I, the first
prime condition makes that vertex good. If the nice point is
exceptional, the original good prime serves as the stable
generization for the length-zero second prime condition. This
again makes the vertex good.

\proofstep{The induction along the component graph.}
Along a path of adjacent vertices keep the following information.
A type-I vertex reached is good. At a thick vertex keep a
nice exceptional point containing it, linked by a finite
sequence of nice exceptional points to one having a good type-I
generization. In the chart ring \(A\), adjacency of successive vertices \(P,P'\) supplies a common specialization of controlled dimension: 
\[
 V\supseteq P+P',\qquad\dim A/V\ge D.
\]
 Choose the nice specialization \(Q\) above \(V\); then
\[
                       P+P'\subseteq V\subseteq Q .
\]

If \(P\) is good type I, use \(M_n\) in the specialization
and apply the first or the length-zero second condition at
\(Q\), according to its type. If \(P_1\) is any type-I generization of this \(Q\), either implication gives 
\[
 P_1\subseteq Q\quad\Longrightarrow\quad M_n\subseteq P_1\cap B_n.
\]
 If \(P\) is thick, \(V\)
contains the corner ideal \(I_3\). Use \(I_3\) in the
specialization: \(Q\) still contains the odd ideal, so is
exceptional.
Write \(Q_{\rm old}\) for the stored nice point on the current thick vertex. The common vertex joins the new exceptional point to the chain since 
\[
 P\subseteq Q_{\rm old}\cap Q.
\]

The second prime condition now makes every type-I generization
of the new point good. If the next vertex is type I, this proves
the assertion for it. If it is thick, specialize \(V\) using
its \(I_3\) as well. A previous vertex of type I gives a good type-I generization of this new exceptional point. When the previous vertex is
thick the stored chain extends through that common vertex.
This proves the invariant without an assumption that every
thick component belongs to the rank-three Hecke locus.
In rank two all the vertices and nice points used are type I.

\proofstep{Passage between charts.}
It remains to pass between charts. Write \(\pi_0\) for the projection to the degree-zero spectrum, \(x_0\) for its residual point and \(Z_{\rm up}\) for the upper triangular locus. The core used to join the charts is 
\[
 C=\pi_0^{-1}(x_0)\setminus Z_{\rm up}.
\]
 Its relevant points admit smooth lower
triangular realizations over their residue fields. Minimal
frame primes are homogeneous. A homogeneous type-I generic
point specializes to this core. One chooses a lower entry nonzero
there, uses the chart obtained by retaining its direction under
the torus, and specializes its degree-zero coefficients to the
closed point.  Applying this to a homogeneous lift of a seed contraction gives a chart prime \(\mathfrak q\) with quotient dimension
\[
 \dim(\text{chart}/\mathfrak q)=\dim\bigl(B_n/(P\cap B_n)\bigr)+1 .
\] The extra parameter is the one-dimensional torus orbit.

The core is preconnected. In rank two the sole lower entry is
a twisted cocycle subject to linear local and marked equations.
Interpolation over \(k[t]\) to a fixed nonzero solution gives
a common generization, still off the upper triangular locus.
In rank three write \(y\) for the adjacent lower entry and
\(a\) for the outer lower entry. The other adjacent entry
is given by the involution.

Write \(\partial a\) for the linear part of the outer
cocycle equation and \(Q(y)\) for its quadratic term.
A pure outer solution \(b\) satisfies \(\partial b=0\),
so
\[
        \partial(t^2 a+(1-t)b)=t^2 Q(y)=Q(ty).
\]
This is the lower multiplication law for the interpolant
\((ty,t^2a+(1-t)b)\). The local flag coordinates scale
in degrees one and two, and adding a pure outer solution adds
its outer flag coordinate. Pick \(b\ne0\) as supplied by the
frame. The kernel of the map into \(k[t]\) is a common
generization of \((0,b)\) and the given point. To see density
off the exceptional locus one needs a solution with \(y\ne0\).

For the frame's nonzero adjacent cocycle \(u\), put \(P=\bar\Psi\)
and use the conventions \(u(gh)=P(h)u(g)+u(h)\) and
\(u(g^*)=u(g)/P(g)\). Set, pointwise,
\[
 y=u+\delta(P-1),\qquad
 a=\frac{u^2}{2P}-\delta u(1+P^{-1})
                +\delta^2\bigl(1-(P+P^{-1})/2\bigr).
\]
Here \(\delta(P(g_*)^2-1)=-u(g_*)(P(g_*)-1)\), with
\(g_*\) the marked element for the normal form. Substitution
gives the marked equality for \(y\), the quadratic multiplication
law for \(a\) and the required local coboundary forms, using
\(u|_{G_v}=\alpha_v(P-1)\). If \(y=0\), substitution in the formula for \(u(g_*^*)\) gives
\[
 2\delta\bigl(P(g_*)^2-1\bigr)=0 .
\]
This forces \(u=0\), a contradiction. With one such solution
\((y_0,a_0)\), the following curve exhibits a pure outer point as a specialization: 
\[
 (y(t),a(t))=(t y_0,t^2a_0+a),\qquad (y(0),a(0))=(0,a).
\]
 Since \(y_0\ne0\), the generic point lies outside the exceptional locus.

\proofstep{From the open core to all required components.}
Divide the finitely many components meeting the open type-I
locus according as their generic points are good or bad.
The high-dimensional seed in its pointed chart gives a good
component by the local propagation. Every component under
consideration meets the core by marked specialization.
The good and bad unions are closed and cover the core, because
off the exceptional locus its closed points lie on type-I
components by the ungraded chart calculation, and these
points are dense there and in the core.

By the Jacobson
specialization property of this fiber, a closed set
meeting it contains a closed chart center there, also after
restriction to an open subset of the core.
Write \(C_{\rm good}\) and \(C_{\rm bad}\) for the two closed subsets of the core \(C\). Preconnectedness gives 
\[
 C=C_{\rm good}\cup C_{\rm bad},\quad C_{\rm good},C_{\rm bad}\ne\varnothing\quad\Longrightarrow\quad C_{\rm good}\cap C_{\rm bad}\ne\varnothing.
\]
 Choose a closed chart center in the
intersection. Flatness lifts both generic points by going
down, giving two type-I minimal chart primes. Local propagation
from the good one reaches the other, contrary to the latter's
contraction. Thus every type-I prime is good. This argument
first applies to \(M_2\), then to each of \(M_3\) and \(M_0\)
on the same rank-three frame.

\end{proof}

\subsection{Two levels and the local annihilator}

Support at the two levels supplies local information at different exceptional places.
To compare the resulting cuspidal attachments, let \(\xi:A_3\to\overline{\Q}_\ell\) be a continuous integral
type-I point with arithmetic ordinary characters on the first
level. At the other level identify the local groups via the
marked place, in the order prescribed by the common frame.
The map of pairs on each level gives the characters of precisely
the same twisted frame point. Thus local algebraicity and the
strict inequalities of its weights transfer at each embedding.

Both determinant points lie on their respective Hecke supports.
Classicality gives cuspidal attachments with the prescribed
local compatibility away from the respective exceptional places.
Both pullbacks give the same representation \(\epsilon^{-1}r_\xi\). Thus for a good Frobenius \(g\) their common characteristic polynomial is
\[
 \det\bigl(1-X\epsilon(g)^{-1}r_\xi(g)\bigr).
\]
Agreement at almost all unramified places identifies the cusps by strong multiplicity one. At
each exceptional prime the other level gives the missing local
comparison.

\paragraph{Weights for fixed-weight adjoint vanishing.}
The comparison at the two levels will be used locally on a component. Fix an integral
type-I point of the full frame ring \(A=A_n\)
with kernel \(q\). Lemma~\ref{lem:ordinary-annihilator-derivative}
will show that \(a\in M_n A\) vanishes in \(A_q\).
To supply the local hypothesis for the density argument on a component,
choose one embedding \(\tau_v:L_v\hookrightarrow\overline{\Q}_\ell\)
for each marked place over \(\ell\). At each marked place and \(j=0,1,2\), prescribe the weight according to the chosen embedding by
\[
 h_j(\tau)=\begin{cases}2j&\tau=\tau_v,\\ j&\tau\ne\tau_v.\end{cases}
\] The indices increase strictly at every
embedding. The local degree is greater than five, so choose
a nonzero trace-zero \(x\in L_v\), sufficiently small that
the exponential and logarithm are inverse on its neighborhood
of zero. Then \(u=\exp(x)\) is a nontorsion unit with norm one.

Take an inertial Artin lift of \(u\). Its cyclotomic value is
one, and the \(j\)-th ordinary character there has value
\[
   \zeta_j\,N_{L_v/\Q_\ell}(u)^{-j}\tau_v(u)^{-j}
      =\zeta_j\tau_v(u)^{-j},\qquad \zeta_j\ \text{of finite order}.
\]
Every ratio for distinct indices has infinite order on this
single element. The argument works for any finite parts in the
specialization. Density at this weight therefore supplies the
local hypothesis of fixed-weight adjoint vanishing at an
integral type-I point on each component. At
rank three the two-level comparison supplies its finite-place
compatibilities. We give the commutative algebra in detail.

\begin{lemma}[Differentiation on the ordinary support]
\label{lem:ordinary-annihilator-derivative}
With the finite weight quotient and type-I density at these
two levels, an ordinary Hecke relation \(a\) restricted to the
chart is annihilated by an element outside the prime of the
type-I point. Writing \(\xi\) for evaluation at the point and \(D_\xi\) for any
point derivation, this gives
\[
                      \xi(a)=0,\qquad D_\xi(a)=0
\]
also at a point with resonant weights.
\end{lemma}

\begin{proof}
\proofstep{Heights in a flat pointed chart.}
Use a flat pointed chart \(R\) containing a prime \(Q\) over
\(q\), minimal in that fiber, with \(\dim(R/Q)=1\). This chart
is obtained by specializing an integral model of the point
and completing at its residual center. Every prime \(P\subset Q\)
contracts to a type-I prime of the frame, so its degree-zero
contraction contains \(M_n\). That quotient is finite over an
ordinary weight quotient and has dimension at most \(D\).
The chart dimension formula adds at most one for the torus
fiber, giving \(\dim(R/P)\le D+1\). Extend any chain ending at \(Q\) by the specialization to the closed point. The resulting dimension bound is
\[
 \operatorname{ht}Q+1\le D+1 .
\]

In the presentation of \(R\) with \(N\) variables over a
DVR and \(N-D\) relations, the lift of \(Q\) has height \(N\)
by catenarity and \(\dim(R/Q)=1\). Its height after imposing
\(N-D\) equations is at least \(D\). Hence equality holds.
At the lifted prime the relations have their full possible
height in a regular local ring and form a regular sequence.
The quotient local ring \(R_Q\) is therefore Cohen--Macaulay.
For the flat local map \(A_q\to R_Q\), the depth and dimension
formulas with the zero-dimensional closed fibre give
\[
 \operatorname{depth}A_q=\operatorname{depth}R_Q
      =\dim R_Q=\dim A_q .
\]
In particular \(A_q\) is Cohen--Macaulay and has no embedded
associated primes. The dimension identity is
\(\operatorname{ht}q=\operatorname{ht}Q\).

\proofstep{Generic reducedness along every component.}
For each minimal prime \(P_0\subset q\) of \(A\), going down
and taking a minimal prime in the chart lifts it to a minimal
chart prime below \(Q\). Every such chart component has
dimension at least \(D+1\), by the presentation and the
height theorem. Its frame contraction has degree-zero quotient
dimension at least \(D\), by the bound of one on the torus
fiber. If \(P\) is the frame contraction and \(B_n\) its degree-zero ring, the fiber bound and the Hecke upper bound now agree: 
\[
 D\le\dim B_n/(P\cap B_n)\le D.
\]
 Apply density on this component and lift using the
graded frame to a type-I point \(\xi'\) on \(P_0\) with
the weights just prescribed.

At \(\xi'\) the ordinary adjoint
vanishing says every fixed-weight deformation is inner.
In the marked basis the remaining infinitesimal change of
frame has dimension at most one. The extra inertial weight directions contribute at most \(nd\), so the bound for the full tangent space is
\[
 \dim T_{\rm full}\le 1+nd=D .
\] The height lower
bound at \(q'=\ker\xi'\) is \(D\) by the same chart
presentation. Over the characteristic-zero coefficient field
the residue field extension is separable and formally smooth,
so every cotangent functional extends to a coefficient-relative
point derivation. Together the tangent and height bounds give 
\[
 D\le\dim A_{q'}\le\operatorname{edim} A_{q'}\le D.
\]
 Hence \(A_{q'}\) is regular local and a domain; localization at its minimal prime gives the field \(A_{P_0}\).

\proofstep{Differentiating the annihilator relation at the original point.}
Now if \(\operatorname{Ann}_A(a)\subseteq q\), take a prime
minimal over that annihilator inside \(q\). It is associated
to the cyclic submodule \(Aa\), hence to \(A\). It must
therefore be a minimal prime \(P_0\) by the unmixedness at
\(q\). Goodness at \(P_0\) puts \(a\) in \(P_0\). In the
field \(A_{P_0}\) it vanishes. This supplies an annihilator
outside \(P_0\), a contradiction. We obtain the asserted
\(s\notin q\). Since \(s\notin q\), we have \(\xi(s)\ne0\). Evaluation first gives \(\xi(a)=0\); differentiation then gives 
\[
 0=D_\xi(sa)=\xi(s)D_\xi(a)+\xi(a)D_\xi(s)=\xi(s)D_\xi(a).
\]
 Thus every point derivation kills \(a\), even if the original point has resonant
weights.  At that point Section~\ref{sec:rank-three-multiplicity}
supplies the selected multiplicity used on the determinant
quotient in the fixed-weight comparison.

\end{proof}

\paragraph{Exclusion of an adjacent cyclotomic ratio.}
We now use the bound on fixed-inertia derivatives to rule out a cyclotomic ratio for
the marked points specified below.

\begin{paramdefinition}[The points tested for an adjacent cyclotomic ratio]\label{hyp:cyclotomic-ratio-point}
On the rank-three type-I chart of
Proposition~\ref{prop:ordinary-component-propagation}, let
\(X_e\), \(e\ge1\), denote the field points having the finite
inertia parts and local degree bounds of
Proposition~\ref{prop:separated-marked-host}, absolutely irreducible
on each open global subgroup, and potentially crystalline at the
marked \(\ell\)-adic places with ordered exponents \((e,0,-e)\).
At such a point fix the inertia characters when taking derivatives,
as in Lemma~\ref{lem:ordinary-annihilator-derivative}.
\end{paramdefinition}
For the application to the even level, the restriction and
attachment arguments proving Lemma~\ref{lem:rank-three-host-restriction}
and Proposition~\ref{prop:ordinary-host-attachment} supply these
properties before applying the ratio test.

\begin{proposition}[Exclusion of cyclotomic resonance]
\label{prop:ordinary-no-cyclotomic-resonance}
At a point of Definition~\ref{hyp:cyclotomic-ratio-point}, writing \(D_v\) for the marked local group, any two distinct ordinary flag positions satisfy 
\[
 (\psi_i/\psi_j)|_{D_v}\ne\epsilon|_{D_v}\qquad(i\ne j).
\]
 The two estimates concern the same
fixed-inertia tangent space: the global calculation bounds its
rank by one, while the proposed ratio forces rank at least two.
\end{proposition}

\begin{proof}
\proofstep{The upper bound on fixed-inertia derivations.}
First consider the upper estimate at the point \(x\) of the
level determinant ring. A derivation induced from the frame
with fixed inertia annihilates \(J\) by the local annihilator
and kills the power series weight ring \(\Lambda\) by its
generators. Let \(J'\) be the reduced support ideal of the
projected module. A trace Gram minor \(b_0\) nonzero at \(x\)
restricts the support to its irreducible locus. On this open,
dense good arithmetic points of the ordinary Hecke component
have positive selected multiplicity. Thus \(b_0a\in J\) for
\(a\in J'\), by reducedness and density in the Hecke quotient.
The local annihilator gives \(D(b_0a)=0\). Since \(a(x)=0\) and \(b_0(x)\ne0\), 
\[
 0=D(b_0a)=b_0(x)D(a)+a(x)D(b_0)=b_0(x)D(a),
\]
 so \(D(a)=0\).

\proofstep{Separating the eigenvalues in the fixed-weight fibre.}
On the finite free projected weight module choose a linear
combination \(t\) of Hecke generators primitive for its
generic reduced algebra and separating the systems in the
weight fiber at \(x\). Finitely many proper linear conditions
on the coefficients suffice for both separations in
characteristic zero. Let \(t\) be this separating combination of Hecke generators. The characteristic polynomial in the generic algebra has the factorization
\[
 P_t(X)=\prod_i g_i(X)^{m_i},\qquad \gcd(g_i,g_j)=1\quad(i\ne j),
\]
with the \(g_i\) distinct monic irreducibles. They lie in \(\Lambda[T]\) by integral closedness.
Let \(k_i\) be the multiplicity of \(x(t)\) as a root of
the specialization of \(g_i\). If \(k_i>0\), lying over in
the corresponding finite integral component gives a point
with these values of \(\Lambda,t\). Separation identifies it
with \(x\).

On that component choose a good arithmetic point
\(y\) with the same algebraic exponents and with
\(b_0g_i'(t)\prod_{j\ne i}g_j(t)\ne0\). This uses density
at those exponents with varying finite-order parts. The
displayed element is generically nonzero there by separability.
The selected joint dimension at \(y\) is at most the root
multiplicity \(m_i\), and at \(x\) it is the whole root
multiplicity since \(t\) separates at \(x\). Consequently
\[
       \sum_i m_i k_i=\mu_e(x)=\mu_e(y)\le m_i
                              \qquad(k_i>0).
\]
Some \(k_i\) is positive. Exactly one is positive and it is
one. The square-free product \(\prod g_i(t)\) belongs to \(J'\)
and has a simple root at \(x(t)\). Differentiating gives
\(D(t)=0\). The component ideal for that root is killed by
\(D\) as well, since multiplying each element by
\(\prod_{j\ne i}g_j(t)\), a unit at \(x\), puts it in \(J'\).

For \(a\) in the component algebra define
\[
          P_a(T)=\operatorname{Tr}\bigl(a\,g_i(T)/(T-t)\bigr)
                     \in\Lambda[T].
\]
Here the trace is from the generic component field, and
integrality of the coefficients follows from monic division
by \(T-t\) and integral closedness. At \(T=t\) the identity
embedding contributes \(a g_i'(t)\) and all other embeddings
contribute zero. Differentiate this identity at \(x\). The coefficients and \(t\) have zero derivative. Differentiating the surviving term gives 
\[
 g_i'(x(t))D(a)=0,\qquad g_i'(x(t))\ne0,
\]
 and hence \(D(a)=0\). Thus the derivation
vanishes on the whole determinant ring and on its image \(B_3\).

At a type-I frame point there are nonzero specialized entries
of both signs and one of odd weight. Products give two
monomials \(p_+,p_-\) of weights \(1,-1\) nonzero there.
A derivation vanishing on \(B_3\) is determined by its value
on \(p_+\). Write \(\xi\) for the frame point above \(x\). If \(D(p_+)=0\), applying \(D\) to \(p_+p_-\) gives \(D(p_-)=0\). For an entry \(u\) of weight \(j\ne0\), take \(p=p_-\) if \(j>0\) and \(p=p_+\) if \(j<0\). Its balanced product then gives 
\[
 0=D(up^{|j|})=\xi(p)^{|j|}D(u),\qquad \xi(p)\ne0.
\]
 All generating entries
and ordinary characters have zero derivative. This proves the
upper bound one for the dimension of fixed-inertia derivations.

\proofstep{The lower bound from a cyclotomic ratio.}
For the lower bound, equality on an open inertia subgroup forces equality of cyclotomic exponents since \(\epsilon\) has infinite image there:
\[
 e(w_i-w_j)=1,\qquad (w_1,w_2,w_3)=(1,0,-1).
\] Thus \(e=1\) and
the two indices are adjacent. Normalize the adjacent
two-dimensional subquotient by its second character. Its
extension cocycle \(b\) takes values in \(E(1)\), with \(E\)
a finite coefficient field. Its class belongs to the
Bloch--Kato finite subspace. Indeed the subquotient and its
characters become crystalline after finite extension and
restriction on cohomology with crystalline period coefficients
is injective by averaging in characteristic zero. By Kummer
theory this finite subspace consists of local units, annihilated
under cup product by the unramified line in \(H^1(F_v,E)\).
For a nonzero unramified additive character \(u\) take a
cochain satisfying
\[
      \beta(st)=\beta(s)+\epsilon(s)\beta(t)-b(s)u(t).
\]
Scaling \(u,\beta\) together makes them integral by compactness.
\cite{bloch-kato-finite}

Write \(A_k=\mathcal O_E/\ell^k\), \(q_k=|A_k|\), and work
in a stable ordinary integral basis. At each \(v\mid\ell\)
the local Euler characteristic and Tate duality give at least
\(q_k^{6(1+d_v)}\) upper triangular adjoint cocycles before
fixing inertia. At each of the \(d_v\) pro-\(\ell\) unit generators, three diagonal coordinates are fixed. Thus the maximum cost is
\[
 (q_k^3)^{d_v}=q_k^{3d_v}.
\]
The marked local fields have no \(\ell\)-power roots of unity.
At the distinguished place impose equal variations on the
adjacent diagonals. The resulting module has rank five and
projects equivariantly onto \(A_k(1)\) via the adjacent entry,
so the dual invariant factor has size at least \(q_k\). Only
two inertia characters must be killed. With the two diagonal variations equal, the lower bound for the number of solutions simplifies to
\[
 q_k^{5(1+d_v)+1-2d_v}=q_k^{6+3d_v}.
\]

One further integral cocycle has unramified diagonal difference
\(u\) on the resonant pair and zero difference on the other
adjacent pair. Its resonant entry is supplied by \(\beta\). Then
the remaining error is a two-cocycle in the outer root
\(\kappa_0/\kappa_2\). This differs from \(\epsilon\)
since the other adjacent ratio is nontrivial at the marked
regular element. With \(\theta=\kappa_0/\kappa_2\), local duality identifies the obstruction space by 
\[
 H^2(G_v,E(\theta))^\vee\simeq H^0(G_v,E(\epsilon\theta^{-1}))=0,
\]
 since \(\epsilon\theta^{-1}\) is nontrivial. Correct the outer entry and clear denominators
once.

Multiples of this integral cocycle at depth \(k\)
supply at least \(q_k/|\mathcal O_E/(c)|\) different cosets
of the equal-diagonal solutions, where \(c\ne0\) is a fixed
value of its diagonal difference. Adding all full-matrix
coboundaries gives another factor at least \(q_k^3\) at
every place, since distinct residual eigenvalues at the
marked local element force the generator of any coboundary
that stays upper triangular to be upper triangular. Thus
the allowed local cocycles \(Z_v\) satisfy
\[
  \prod_{v\mid\ell}|Z_v|
        \ \ge\ q_k^{9\,\#\{v\mid\ell\}+3d+1}/C
\]
with \(C\) independent of \(k\) for this point.

\proofstep{The global cardinality estimate and contradiction.}
The global Selmer cardinality inequality for these local
conditions gives
\[
 |Z_{\rm glob}|\, q_k^{9\,\#\{v\mid\ell\}}
             \prod_{v\mid\infty}|H^0_v|
      \ \ge\ q_k^9\prod_{v\mid\ell}|Z_v| .
\]
The dual global invariant factor is trivial here. The
residual frame separates each matrix unit's adjoint character
from \(\bar\epsilon\), using \(1,\bar\Psi^{\pm1}\) on the
positive component and \(\eta\bar\Psi^{\pm2}\) on the outer
roots. Nakayama shows that the images of
\(\operatorname{Ad}\rho(s)-\epsilon(s)\) span the truncated
module, so no equivariant dual functional survives. At each real place the fixed matrices identify with the skew matrices; there are three free entries.

Taking the product over the \(d\) real places gives
\[
 \prod_{v\mid\infty}|H_v^0|=(q_k^3)^d=q_k^{3d}.
\] Before normalization the global cocycle bound is \(q_k^{10}/C\). Imposing zero
at conjugation costs six scalar entries by its square relation.
For \(Y=z(g_1)\rho(g_1)\) the remaining equations set the
off-diagonal part of \(Y-JY^tJ\) to zero, with \(J\) the
antidiagonal matrix. The two outer corners already vanish
and the other four entries form two opposite pairs. Thus, writing \(Z_k\) and \(Z_k^{\rm norm}\) for the cocycles before and after imposing the two normalizations, respectively, we obtain
\[
 |Z_k^{\rm norm}|\ge q_k^{-6-2}|Z_k|
              \ge q_k^{-8}q_k^{10}/C=q_k^2/C .
\]

Each normalized cocycle deforms the representation by
\((1+\delta z(s))\rho(s)\), \(\delta^2=0\), with the
prescribed ordinary conditions, and hence induces a distinct
fixed-inertia point derivation into \(A_k\). The relative
cotangent module after imposing those inertia equations is
finite over \(\mathcal O_E\). Write \(M\) for this module and \(T\) for its torsion submodule. If \(\operatorname{rk}M\le1\), the elementary divisor decomposition over the coefficient DVR gives
\[
  q_k^2/C\le |\operatorname{Hom}_{\mathcal O_E}(M,A_k)|
      \le q_k |T| .
\]
Here \(C\) and \(T\) are fixed as \(k\) varies. Since \(q_k\to\infty\), the rank is at least two. This contradicts
the upper bound on the very same framed tangent space and
excludes the cyclotomic equality. The constants in this
depth argument can depend on the chosen point and the field.

At a type-I point of the even-level Field used in the rank-three trace estimate over the
imaginary quadratic field \(K\), the marking for this use takes
\(M=\ell N|\operatorname{disc}K|\). On complex conjugation \(c\) the two parity factors cancel:
\[
 (\epsilon^e\psi)(c)=(-1)^e\psi(-1)=(-1)^{2e}=1 .
\] Every prime ramifying in the quadratic
character divides \(M\). Factor the integral point through
the integers of a finite coefficient extension. Its ordinary
flag pulls back to the marked groups and its restriction to
the open CM subgroups is irreducible by the type-I host
hypothesis.

Given the automorphic attachment on this field, the preceding result supplies the unlinked local flag at the even level. The resulting local nonresonance and the three distinct weights give the required freedom of the ordinary characters. The original packet level in this branch
is \(N\), while \(M\) is the ramification cutoff of the comparison.
For the auxiliary field over \(\mathbb Q(i)\), use the same local exclusion after lifting to that field in the rigidity and trace bounds. Thus the local calculation is shared
even though the host attachment and the two-dimensional
constituent of the other Case of the rank-three trace estimate require separate
arguments.
\end{proof}
